% Benchmark suite: algorithms, implementations, and swept parameters.
% Self-contained section; uses only booktabs (for \toprule/\midrule/\bottomrule).
% \usepackage{booktabs}

\section{Benchmarks}
\label{sec:benchmarks}

\subsection{Methodology}
\label{sec:methodology}

Every benchmark is a first-order recurrence whose per-step state can be
\emph{folded} (stored in a small sliding window) rather than \emph{materialized}
over the whole recurrence axis. To separate the effect of storage folding from
the effect of producer/consumer fusion, each application is expressed as three
variants that share the same algorithm and fusion structure and differ only in
how much of the recurrence state is kept live:

\begin{enumerate}
  \item \textbf{non-inductive (materialize)} --- the baseline formulation, written
    \emph{without} inductive functions: the full trajectory (or a dense sliding
    window via an \texttt{RDom}) is materialized in a buffer. This is the
    ``obvious'' Halide implementation.
  \item \textbf{inductive unfolded} --- the inductive-function formulation with
    storage folding \emph{disabled} (fold factor pinned to the full
    recurrence-axis extent). It differs from (3) only in the fold factor, so it
    isolates \emph{folding} from \emph{fusion}.
  \item \textbf{inductive folded} --- the inductive-function formulation with the
    recurrence axis folded to the minimal window (typically $2$--$3$ slots). This
    is the proposed implementation.
\end{enumerate}

Where a mature third-party or hand-written baseline exists (OpenCV for stereo,
Boost.odeint for the ODE, oneTBB/CUB/Thrust for scan, the vendor CUDA kernel for
Mamba) it is included as an \emph{external reference}; it is not a non-inductive
\emph{formulation} in the sense above and is reported separately.

\paragraph{Why these computations are inductive.}
Each benchmark is a \emph{recursively-defined producer feeding a consumer}: a Func
whose update reads its own value at a bounded lag along one axis, whose result is
then consumed downstream (a reduction, a running mean, an output image, an
observer). Written naively, such a producer must be computed and stored in full
before its consumer runs (\texttt{compute\_root}), which both destroys
producer/consumer locality and materializes the entire trajectory. An inductive
function instead lets the recursive producer be computed \emph{fused with} its
consumer---each recurrence step is produced exactly where it is consumed, so the
two remain resident together in cache or registers---while the recurrence axis is
kept in a small \emph{folded} window rather than in full. Two properties make this
legal, and every benchmark below has them: \textbf{(i)~bounded-lag self-reference}
(a $k$th-order recurrence keeps only $k{+}1$ slots live), and
\textbf{(ii)~monotone in-order consumption} (the consumer reads each step before
folding overwrites it). The payoff is therefore twofold---preserved
producer/consumer locality \emph{and} storage folding of the recurrence state from
the whole trajectory down to the window---and it is largest exactly when the
unfolded trajectory is big relative to cache (all apps except Chebyshev, whose
trajectory stays cache-resident and is included as the in-cache tie).

\paragraph{Dimension roles.}
For each application we distinguish three kinds of axis, because folding interacts
with each differently:
\begin{itemize}
  \item \textbf{recurrence length} --- the axis the recurrence runs along, i.e.\ the
    \emph{folded} axis. Sweeping it grows the unfolded footprint without changing
    per-step work; it is the axis that crosses the cache boundary.
  \item \textbf{per-step work} --- the computation per recurrence step (transition
    scan width, matvec size, stencil width). Sweeping it changes arithmetic
    intensity.
  \item \textbf{batch / parallel} --- independent problem instances scheduled in
    parallel. Sweeping it grows total footprint via a \emph{parallel} axis, so its
    points collapse onto the same footprint/cache curve as the recurrence-length
    sweep.
\end{itemize}

\paragraph{Measurement protocol.}
Timing is the median of $30$ timed iterations after $3$ warmups (best/min and IQR
also recorded); each iteration is one full \texttt{realize()} (plus a device sync
on GPU), and runs are pinned to a single NUMA node. The recurrence-state
footprint is reported \emph{byte-exact}: either measured as the high-water mark of
a custom Halide allocator wrapped around one untimed \texttt{realize()} (Viterbi,
Chebyshev, ODE) or computed analytically from the schedule (Kalman, prefix-sum,
StereoBM). Input/output buffers are excluded, so only the recurrence state---the
quantity folding controls---is counted. Footprint is deterministic (independent of
CPU speed and, for these applications, of thread count), so it is captured with a
single trial. CPU results use an AMD EPYC (Zen~2, AVX2) node; GPU results an
NVIDIA A10G (sm\_86). All Halide variants use the same Halide build.

\subsection{CPU benchmarks}

\paragraph{Viterbi decoding.}
Maximum-likelihood state-sequence decoding for a hidden Markov model in the log
domain. At each time step the per-state score is updated by a max-plus transition
scan over the previous step's scores plus the emission term, with a backpointer
retained for the traceback; the forward pass is a first-order recurrence over time
$t$.
\emph{Inductive because:} the score column is a recursive producer reading only its
$t{-}1$ predecessor (lag~1, Markov), consumed by the next step (and the traceback)
in increasing~$t$; fusing it with that consumer folds the score axis $t\!\to\!2$.
\emph{Implementations:} non-inductive (materialize the full $S\times T$
score/backpointer trajectory); inductive unfolded ($\mathrm{fold}\;t\!\to\!T{+}1$);
inductive folded ($\mathrm{fold}\;t\!\to\!2$).
\emph{Setups} (Fig.~\ref{fig:combined-setups}):
\begin{itemize}\itemsep2pt
  \item \textit{rec-len (1c)} --- recurrence length swept over
    $T\in\{5000,20000,80000,320000\}$ (4 values, spanning the LLC boundary), with
    $S{=}16$ states and $M{=}4$ observations fixed, single-threaded (Viterbi has no
    parallel axis).
  \item \textit{per-step-work} --- state count swept over $S\in\{4,16,64,256\}$
    (4 values, growing the $O(S)$ transition scan and the footprint together),
    with $M{=}8$ and $T{=}50000$ fixed, single-threaded.
\end{itemize}

\paragraph{Kalman / latent AR(2) log-likelihood.}
Log-likelihood of a latent second-order autoregressive process with observation
noise, evaluated by the Kalman filter recursion (predict/update) over time; the
filtered mean/covariance form a first-order recurrence over $t$ and the
log-likelihood accumulates along the way. Batched over $B$ independent series.
\emph{Inductive because:} the filtered state is a recursive producer needing only
its $t{-}1$ value (the AR(2)'s two lags are absorbed into the augmented state, so
still lag~1), consumed by the running log-likelihood reduction in time order; fused
with that reduction it folds $t\!\to\!2$.
\emph{Implementations:} non-inductive (materialize the filtered trajectory);
inductive unfolded ($\mathrm{fold}\;t\!\to\!T$); inductive folded
($\mathrm{fold}\;t\!\to\!2$).
\emph{Setups} (Fig.~\ref{fig:combined-setups}):
\begin{itemize}\itemsep2pt
  \item \textit{rec-len (1c)} --- recurrence length swept over
    $T\in\{1024,4096,16384,65536,131072\}$ (5 values, crossing the LLC boundary),
    $B{=}256$ series, single-threaded.
  \item \textit{rec-len (32c)} --- the same $T$ sweep (5 values), $B{=}256$, at
    $32$ threads; separated from the 1-core run because core count is a
    qualitatively different regime.
  \item \textit{batch} --- batch (parallel axis) swept over
    $B\in\{16,64,256,1024\}$ (4 values, footprint crosses the LLC via $B$),
    $T{=}16384$ fixed, threads capped to the task count $\min(\mathrm{nproc},B/8)$.
  \item \textit{arith-intensity} --- a single point ($B{=}256$, $T{=}16384$, $32$
    threads) with the per-step divide/log replaced by a fixed-reciprocal multiply
    (same footprint), isolating whether a folding loss is latency- rather than
    footprint-bound.
\end{itemize}

\paragraph{Chebyshev semi-iteration.}
Chebyshev semi-iterative solver for a linear system: each iteration is a dense
matrix--vector product plus the Chebyshev three-term recurrence combining the
current and two previous iterates. The recurrence length is the iteration count
$M$; per-step work is the $O(n^2)$ matvec.
\emph{Inductive because:} each iterate is a recursive producer of the two previous
iterates (lag~2), consumed by the next matvec in increasing~$k$; fused, the iterate
axis folds $\to 3$.
\emph{Implementations:} non-inductive full materialize ($M{+}1$ iterate columns);
non-inductive mod-3 ring ($3$ columns, hand-written modulo indexing); inductive
unfolded ($\mathrm{fold}\!\to\!M{+}1$); inductive folded ($\mathrm{fold}\!\to\!3$).
\emph{Setups} (Fig.~\ref{fig:combined-setups}):
\begin{itemize}\itemsep2pt
  \item \textit{per-step-work} --- system size swept over
    $n\in\{512,1024,2048,3072\}$ (4 values, growing the $O(n^2)$ matvec per
    iteration), with $M{=}100$ iterations fixed, single-threaded.
  \item \textit{rec-len} --- iteration count / fold ratio swept over
    $M\in\{50,100,200,400\}$ (4 values), with $n{=}2048$ fixed, single-threaded.
\end{itemize}
The footprint ($n\,M\cdot 8$) stays cache-resident at realistic sizes, so this is
the intended \emph{in-cache} datapoint (fold $\approx$ tie).

\paragraph{Allen--Cahn ODE integration.}
Method-of-lines integration of the Allen--Cahn reaction--diffusion PDE with a
two-step Adams--Bashforth (AB2) scheme. The spatial field of width $D$ is advanced
by a sparse stencil; AB2 makes the update a first-order recurrence needing the two
previous derivative evaluations, and a fused observer reduces the trajectory on
the fly. Batched over $B$.
\emph{Inductive because:} the field is a recursive producer needing the two prior
derivative slabs (lag~2, AB2), consumed immediately by the fused observer in time
order; that fusion is what folds time $t\!\to\!3$ and avoids storing the
space$\times$time trajectory.
\emph{Implementations:} non-inductive (materialize the space$\times$time
trajectory); inductive unfolded ($\mathrm{fold}\;t\!\to\!T{+}1$); inductive folded
($\mathrm{fold}\;t\!\to\!3$, the AB2 window). External references: Boost.odeint (RK
integrator) and a C++ reference$+$observer oracle.
\emph{Setups} (Fig.~\ref{fig:combined-setups}):
\begin{itemize}\itemsep2pt
  \item \textit{rec-len} --- number of time steps swept over
    $T\in\{2048,8192,32768,131072\}$ (4 values, crossing the LLC boundary), with
    spatial width $D{=}1024$ and batch $B{=}1$ fixed, single-threaded.
  \item \textit{batch} --- batch swept over $B\in\{1,4,16\}$ (3 values, growing
    the footprint $D\,B\,T\cdot 4$), with $D{=}1024$ and $T{=}8192$ fixed,
    single-threaded.
\end{itemize}

\paragraph{Batched prefix scan.}
Inclusive prefix sum over each row of a $W\times H$ matrix followed by a
running-mean (or cheap $\gg 2$) consumer. The scan is a first-order additive
recurrence along $W$; associativity of $+$ also admits a two-stage chunk
decomposition that parallelizes \emph{along} the scan axis.
\emph{Inductive because:} the prefix is a recursive producer needing only its
$i{-}1$ value (lag~1, a single carried scalar), consumed left-to-right by the
running-mean/shift stage; fused, the scan axis folds to one accumulator
($x\!\to\!1$).
\emph{Implementations:} inductive folded ($\mathrm{fold}\;x\!\to\!1$, $O(1)$
accumulator); inductive unfolded ($\mathrm{fold}\;x\!\to\!W$, materialized prefix
row); non-inductive (\texttt{RDom}, materialize the $O(W)$ row); one-stage folded
(parallel over rows only); two-stage inductive scan (parallel along time via
chunking). External reference: oneTBB \texttt{parallel\_scan} (and, on GPU,
CUB/Thrust---reported separately as vendor scans).
\emph{Setups} (Fig.~\ref{fig:combined-setups}) --- all on a fixed
$W{=}1048576 \times H{=}32$ matrix (32 rows of length $\sim\!10^6$); the two
qualitative axes are the downstream consumer and the core count, and each cell is
a single operating point (no parameter is swept within it, hence no whisker):
\begin{itemize}\itemsep2pt
  \item \textit{mean (1c)} / \textit{mean (32c)} --- running-mean consumer
    ($p_i/(i{+}1)$), single-threaded and at capped threads ($\min(\mathrm{nproc},32)$).
  \item \textit{shift (1c)} / \textit{shift (32c)} --- the cheap $\gg2$ consumer
    (isolates the scan from the divide), single-threaded and at capped threads.
\end{itemize}
(A separate parallelism-regime sweep over lane count $S\in\{1,8,32\}$ ---
one-stage vs.\ oneTBB vs.\ two-stage, chunk $L{=}4096$ --- has no non-inductive
counterpart and so does not appear in Fig.~\ref{fig:combined-setups}.)

\paragraph{Stereo block matching.}
StereoBM disparity estimation: for each pixel and candidate disparity, a
sum-of-absolute-differences over a $\mathit{winsize}\times\mathit{winsize}$ support
window, selecting the disparity that minimizes SAD. The vertical accumulation of
the SAD box filter is a first-order sliding-window recurrence down image rows $y$,
which is what folds.
\emph{Inductive because:} the column running-sum is a recursive producer of its
$y{-}1$ value (lag~1: add the entering row, subtract the leaving one), consumed
top-to-bottom by the disparity-selection stage; fused, the vertical accumulation
folds $y\!\to\!1$.
\emph{Implementations:} inductive folded ($\mathrm{fold}\;y\!\to\!1$, one running
row); inductive unfolded ($\mathrm{fold}\;y\!\to\!H$, full column); schedule4
(non-inductive, \texttt{RDom} vertical slide). External reference: OpenCV StereoBM
on the same image pair.
All setups run on the real Middlebury \emph{Aloe} stereo pair and, within each,
sweep the strip width $\mathit{tilesize}\in\{8,12,16,24,32,40,48,64,80,96,128\}$
(11 values) --- a pure scheduling knob that does not change the result (so OpenCV
parity holds) but trades cache locality against the size of the materialized
non-inductive/unfolded column. Because $\mathit{tilesize}$ affects only
performance and not the result (unlike the other apps' swept parameters, which are
different problems), each stereo setup in Fig.~\ref{fig:combined-setups} reports
the \emph{best tilesize per implementation} --- the best achievable folded/unfolded
time against the best achievable non-inductive time --- rather than a geomean over
the sweep, and so appears as a single point with no whisker. The qualitative axes
are image size (which sets the core regime) and SAD window size:
\begin{itemize}\itemsep2pt
  \item \textit{small w9} --- the small Aloe pair \texttt{aloeL/R.png}
    ($307\times265$), $\mathit{winsize}{=}9$, \textbf{1 core} (single-threaded).
  \item \textit{full-res w9} --- the full-resolution pair
    \texttt{aloeL/R\_large.png} ($1282\times1110$), $\mathit{winsize}{=}9$,
    \textbf{multi-core}: threads are capped to the tile count
    $\min(\mathrm{nproc},\lceil W/\mathit{tilesize}\rceil)$, so the core count
    co-varies with tilesize --- $11$ to $48$ across the sweep on the $48$-core
    measurement node (which is why image size, not cores, is the setup axis here);
    the best folded and best non-inductive times fall at different tilesizes ($21$
    and $41$ cores respectively).
  \item \textit{full-res w15} --- the same full-resolution pair at the larger
    $\mathit{winsize}{=}15$ window, \textbf{1 core} (single-threaded).
\end{itemize}

\subsection{GPU benchmark: batched prefix scan}

The same batched prefix scan (the CPU \emph{Batched prefix scan} benchmark above)
also runs on the GPU, where the
question is not fold-vs-materialize but whether an inductive Halide scan can hold
its own against hand-tuned vendor scan libraries. Because $+$ is associative, the
scan parallelizes \emph{along} time via the two-stage chunk decomposition
$t=kL+j$: an up-sweep computes each chunk's total (parallel over chunks and
lanes), a short serial $O(C)$ pass forms the inter-chunk carries, and a down-sweep
adds the carry to a per-chunk inductive scan whose $j$-axis folds to a single
register. Chunks map to GPU blocks and lanes to threads.
\emph{Inductive because:} the per-chunk local scan is a recursive producer of its
$j{-}1$ value (lag~1), consumed in order by the down-sweep; fused, the $j$-axis is
never materialized ($\mathrm{fold}\;j\!\to\!1$).
\emph{Implementations:} the Halide inductive two-stage scan, versus two NVIDIA
vendor scans --- Thrust \texttt{inclusive\_scan\_by\_key} (the like-for-like
\emph{segmented} scan) and CUB \texttt{DeviceScan::InclusiveSum} (a single
\emph{unsegmented} array, an easier problem shown only as a bandwidth roofline).
The serial-time schedule forced on a \emph{non}-associative recurrence
(\texttt{prefixsum\_gpu\_nonassoc}, no scan-axis parallelism) exhausts device
resources at these sizes and does not complete, so---as with Mamba---there is no
non-inductive GPU baseline; this is an inductive-Halide-vs.-vendor comparison.
\emph{Swept:} recurrence length via total size $W\!\times\!S \in
\{67\mathrm{M}, 268\mathrm{M}, 1.07\mathrm{B}\}$ elements at $S{=}512$ lanes and
chunk $L{=}4096$ (Fig.~\ref{fig:gpu-prefix}); an occupancy sweep over lane count
$S$ at fixed $W$; and a chunk-length sweep $L\in\{1024,4096,16384\}$ (throughput
varies ${<}1\%$, i.e.\ the fold depth is not throughput-critical once the device
is filled). As Fig.~\ref{fig:gpu-prefix} shows, the inductive scan matches Thrust's
segmented scan to within ${\sim}5\%$ at the sizes where the launch is amortized,
and it is the only implementation that still runs at $1.07$B elements, where both
vendor scans run out of device memory.

% GPU batched prefix-scan throughput: the Halide inductive two-stage scan vs.
% NVIDIA vendor scans, on an A10G (sm_86). Thrust inclusive_scan_by_key is the
% like-for-like SEGMENTED competitor; CUB DeviceScan::InclusiveSum is a single
% unsegmented array -- a hardware-bandwidth reference, not a segmented baseline.
% Recurrence-length sweep at S=512 lanes, chunk L=4096; label = total elements
% W*S. At the largest size both vendor scans OOM (nan) while the folded Halide
% scan, which never materializes the state, still runs.
% Data sidecar: gpu_prefix.dat.
% Preamble needs: \usepackage{pgfplots}\usepackage{pgfplotstable}\pgfplotsset{compat=1.16}

\definecolor{cHal}{HTML}{4C72B0}
\definecolor{cThr}{HTML}{55A868}
\definecolor{cCub}{HTML}{C0C0C0}

\pgfplotstableread[col sep=tab]{gpu_prefix.dat}\gputable

\begin{figure}[t]
  \centering
  \begin{tikzpicture}
    \begin{axis}[
      ybar, bar width=9pt,
      width=8.6cm, height=6.2cm,
      symbolic x coords={67M,268M,1.07B},
      xtick=data,
      enlarge x limits=0.35, clip=false,
      ymin=0, ymax=66,
      ylabel={throughput (Gelem/s)},
      ylabel style={font=\footnotesize},
      xlabel={problem size (total elements $W\!\times\!S$, $S{=}512$ lanes)},
      xlabel style={font=\footnotesize},
      x tick label style={font=\footnotesize},
      y tick label style={font=\footnotesize},
      ymajorgrids, grid style={draw=black!10, line width=0.4pt},
      axis lines=left, tick align=outside,
      legend style={at={(0.5,1.02)}, anchor=south, legend columns=-1, font=\scriptsize,
                    draw=black!25, /tikz/every even column/.append style={column sep=5pt}},
      legend cell align=left, area legend,
    ]
      \addplot+[ybar, fill=cHal, draw=white, line width=0.3pt]
        table[x=label, y=halide]{\gputable};
      \addlegendentry{Halide inductive 2-stage}
      \addplot+[ybar, fill=cThr, draw=white, line width=0.3pt]
        table[x=label, y=thrust]{\gputable};
      \addlegendentry{Thrust segmented}
      \addplot+[ybar, fill=cCub, draw=black!25, line width=0.3pt]
        table[x=label, y=cub]{\gputable};
      \addlegendentry{CUB unsegmented (ref)}
    \end{axis}
  \end{tikzpicture}
  \caption{GPU batched inclusive prefix scan on an A10G (sm\_86). The Halide
    inductive two-stage scan (blue) matches NVIDIA's Thrust segmented scan (green)
    to within ${\sim}5\%$ once the problem is large enough to amortize launch and
    fill the device, and it still runs at the largest size where both vendor scans
    exhaust device memory. CUB's \emph{unsegmented} single-array scan (grey) is
    shown only as a bandwidth reference; it solves an easier problem (no per-lane
    segment boundaries) and marks the roofline rather than a like-for-like
    baseline.}
  \label{fig:gpu-prefix}
\end{figure}


\subsection{GPU benchmark: Mamba/S6 selective scan}

The Mamba selective-scan step is a first-order \emph{affine} recurrence per
(batch, channel, state), $h_t = a_t h_{t-1} + b_t$, with $a=\exp(\Delta A)$,
$b=\Delta B u$, $\Delta=\mathrm{softplus}(\delta)$; the output is a pointwise
reduction $y_t = D u_t + \sum_n C_n h_{n,t}$. Only $h$ is a true recurrence.
\emph{Inductive because:} $h$ is a recursive producer of its $t{-}1$ value (lag~1,
affine), consumed once per step in time order by the output reduction; fused, $h$
is folded to a single register and never materialized.
\emph{Implementations:} a Halide register-fold inductive scan (serial in time,
$h$ kept in registers); a Halide two-stage time-parallel scan (composing affine
maps, $(a_2,b_2)\circ(a_1,b_1)=(a_2 a_1,\,a_2 b_1{+}b_2)$, over chunks); and the
vendor CUDA v1 kernel (\texttt{selective\_scan}). There is \emph{no} non-inductive
Mamba baseline---materializing all of $h$ is ${\approx}8$\,GiB and exceeds the
buffer-size limit---so Mamba is an inductive-Halide-vs.-vendor comparison rather
than an inductive-vs.-non-inductive one.
\emph{Swept} (realistic Mamba-1 dims $D{=}2048,\,N{=}16$): recurrence length
$T\in\{4096,8192,16384,65536\}$ at $B{=}8$; batch/occupancy $B\in\{1,4,8,16\}$ at
$T{=}8192$; and, for the two-stage kernel, the chunk length $L$ (fold depth vs.\
time-parallelism trade-off).

\begin{table}[t]
  \centering
  \small
  \begin{tabular}{@{}llll@{}}
    \toprule
    Benchmark & Recurrence & Folded axis & Non-inductive baseline \\
    \midrule
    Viterbi     & max-plus HMM forward   & time $t\!\to\!2$      & materialize $S\times T$ \\
    Kalman AR(2)& linear filter          & time $t\!\to\!2$      & materialize trajectory \\
    Chebyshev   & 3-term iteration       & iterate $\to 3$       & full / mod-3 ring \\
    Allen--Cahn & AB2 multistep ODE      & time $t\!\to\!3$      & materialize space$\times$time \\
    Prefix scan & additive scan          & index $x\!\to\!1$     & \texttt{RDom} materialize row \\
    StereoBM    & SAD box filter         & row $y\!\to\!1$       & \texttt{RDom} vertical slide \\
    \midrule
    Mamba/S6    & affine (linear) scan   & time (register)      & none (${\approx}8$\,GiB) \\
    \bottomrule
  \end{tabular}
  \caption{Benchmark recurrences, the axis storage-folding collapses, and the
    non-inductive formulation each is compared against. Mamba has no feasible
    non-inductive baseline and is compared to a vendor CUDA kernel instead.}
  \label{tab:benchmarks}
\end{table}
